% 依赖项（必须放在主文档导言区，即 \begin{document} 之前）：
% \usepackage{tabularray}
% \UseTblrLibrary{varwidth}
% \usepackage{enumitem}
%
% longtblr 模板：占满行宽，单元格左对齐并垂直居中。
\begingroup
\nopagebreak[4]
\enlargethispage{1cm}
\footnotesize
\centering

% 中文续表文字；作用域由 \begingroup ... \endgroup 限定在本表内。
\DefTblrTemplate{conthead-text}{default}{（续表）}
\DefTblrTemplate{contfoot-text}{default}{续下页}

\begin{longtblr}[
  caption = {典型AI智能体安全事件汇总},
  label   = {tab:agent_security_incident},
]{
  width   = \linewidth,
  colspec = {
    |Q[wd=1.5cm,l,m]
    |Q[wd=2cm,l,m]
    |Q[wd=1.5cm,l,m]
    |X[1.15,l,m]
    |X[1,l,m]|
  },
  cells   = {l,m},
  colsep  = 3pt,
  row{1}  = {font=\bfseries},
  rowhead = 1,
  % 每行内容与上下边框之间各保留少量空间；由该行最高单元格决定整行高度。
  rowsep  = 2.5pt,
  measure = vbox,
  stretch = -1,
}
\hline
发生时间（持续时长） & 涉事主体 & 事件类型 & 攻击/异常手法 & 核心损失与影响 \\
\hline
2026.01.29--2026.02.03\newline （披露修复窗口）
& OpenClaw / Moltbook平台
& 一键远程代码执行、后端数据库配置泄露
& \begin{itemize}[leftmargin=*,topsep=0pt,partopsep=0pt,itemsep=1pt,parsep=0pt]
    \item 公开报告显示Moltbook后端数据库缺少行级安全策略，暴露约150万条智能体API令牌与3.5万个邮箱地址~\cite{openclawMoltbookBreach2026}。
    \item 官方安全公告确认WebSocket来源校验缺失，漏洞编号CVE-2026-25253（CVSS 8.8）~\cite{openclawCVE2026}。
    \item 安全研究机构披露官方技能市场ClawHub被植入341个恶意技能，用于分发窃密恶意软件~\cite{openclawClawHavoc2026}。
\end{itemize}
& \begin{itemize}[leftmargin=*,topsep=0pt,partopsep=0pt,itemsep=1pt,parsep=0pt]
    \item 第三方安全扫描确认4万至13万余台OpenClaw实例暴露在公网，多数未启用身份验证~\cite{openclawExposure2026}。
    \item 厂商已发布补丁修复相关漏洞~\cite{openclawClawHavoc2026}。
    \item 已泄露的令牌和已植入的恶意技能不会因补丁发布而自动失效或清除~\cite{openclawClawHavoc2026}。
\end{itemize} \\
% \cline{2-5}
\hline 
2025.12.04披露\newline （未见实际利用报告）
& LangChain（langchain-core框架）
& 序列化注入漏洞、密钥泄露风险
& \begin{itemize}[leftmargin=*,topsep=0pt,partopsep=0pt,itemsep=1pt,parsep=0pt]
    \item 官方安全公告确认dumps()/dumpd()函数存在序列化转义缺陷，漏洞编号CVE-2025-68664（CVSS 9.3）~\cite{langgrinchCVE2025}。
    \item 公开报告显示攻击者可通过提示注入影响LLM输出字段，触发反序列化时的非预期对象实例化~\cite{langgrinchCVE2025}。
    \item 官方说明该缺陷在特定配置下可导致环境变量密钥泄露~\cite{langgrinchCVE2025}。
\end{itemize}
& \begin{itemize}[leftmargin=*,topsep=0pt,partopsep=0pt,itemsep=1pt,parsep=0pt]
    \item 官方通报截至披露时未发现该漏洞被实际利用或已造成大规模攻击的确认案例~\cite{langgrinchCVE2025}。
    \item 受影响的核心软件包全球下载量超8亿次，潜在影响面广~\cite{langgrinchCVE2025}。
    \item 厂商已发布补丁版本0.3.81与1.2.5~\cite{langgrinchCVE2025}。
\end{itemize} \\
% \cline{2-5}
\hline 

2025.10.22披露--2026.01.23修复
& Perplexity Comet浏览器智能体
& 零点击间接提示注入、本地文件与凭证泄露风险
& \begin{itemize}[leftmargin=*,topsep=0pt,partopsep=0pt,itemsep=1pt,parsep=0pt]
    \item 安全研究机构披露恶意日历邀请函可内嵌隐藏指令，用户请求Agent处理邀请即可触发~\cite{perplexityCometVuln2026}。
    \item 研究人员演示Agent可被诱导访问本地文件系统并将内容传出至攻击者控制的服务器~\cite{perplexityCometVuln2026}。
    \item 报告说明相关利用路径还可用于操纵密码管理器交互，存在凭证泄露风险~\cite{perplexityCometVuln2026}。
\end{itemize}
& \begin{itemize}[leftmargin=*,topsep=0pt,partopsep=0pt,itemsep=1pt,parsep=0pt]
    \item 该漏洞由安全研究机构Zenity Labs负责任披露，未见公开报告显示已被实际大规模利用~\cite{perplexityCometVuln2026}。
    \item Perplexity于2026年1月确认修复~\cite{perplexityCometVuln2026}。
    \item 事件被研究者用于说明智能体浏览器场景下传统同源策略等Web安全机制可能失效~\cite{perplexityCometVuln2026}。
\end{itemize} \\
% \cline{2-5}
\hline 

2025年（8月披露）\newline （长期持续）
& 美国多家财富500强科技公司（受骗雇主）
& AI辅助虚假身份就业欺诈
& \begin{itemize}[leftmargin=*,topsep=0pt,partopsep=0pt,itemsep=1pt,parsep=0pt]
    \item 官方报告说明相关人员使用Claude构建具有可信专业背景的虚假身份材料~\cite{anthropicThreatIntel2025}。
    \item 报告显示AI被用于完成求职技术测评并交付实际技术工作~\cite{anthropicThreatIntel2025}。
    \item Anthropic说明相关操作者本身技术能力与英语沟通能力有限，高度依赖AI完成日常任务~\cite{anthropicThreatIntel2025}。
\end{itemize}
& \begin{itemize}[leftmargin=*,topsep=0pt,partopsep=0pt,itemsep=1pt,parsep=0pt]
    \item 官方报告说明该操作涉及美国多家财富500强科技公司，骗取的薪资收入被用于规避国际制裁~\cite{anthropicThreatIntel2025}。
    \item Anthropic评估AI已消除此前该类骗局对长期专业培训的依赖瓶颈~\cite{anthropicThreatIntel2025}。
    \item Anthropic识别后封禁相关账号~\cite{anthropicThreatIntel2025}。
\end{itemize} \\
% \cline{2-5}
\hline 

2025年（8月披露）\newline （至少17家机构）
& 至少17家机构（医院、应急机构、政府部门、宗教组织）
& 大规模数据勒索与恐吓
& \begin{itemize}[leftmargin=*,topsep=0pt,partopsep=0pt,itemsep=1pt,parsep=0pt]
    \item 官方报告说明攻击者使用Claude Code自动化执行系统扫描与凭证收集等操作~\cite{anthropicThreatIntel2025}。
    \item 报告显示攻击者利用AI分析被窃数据价值，并协助生成勒索信内容~\cite{anthropicThreatIntel2025}。
    \item 官方说明该活动未采用传统加密式勒索软件，而是以泄露数据作为威胁手段~\cite{anthropicThreatIntel2025}。
\end{itemize}
& \begin{itemize}[leftmargin=*,topsep=0pt,partopsep=0pt,itemsep=1pt,parsep=0pt]
    \item Anthropic披露该活动涉及至少17家机构，部分案例勒索金额超过50万美元~\cite{anthropicThreatIntel2025}。
    \item Anthropic识别相关活动后封禁涉事账号，并向有关部门通报~\cite{anthropicThreatIntel2025}。
    \item 官方报告将此列为AI被用于自动化执行网络犯罪操作性环节（而非仅提供建议）的案例~\cite{anthropicThreatIntel2025}。
\end{itemize} \\
% \cline{2-5}
\hline 

2025年中\newline （约30个目标机构）
& 约30家全球机构（大型科技公司、金融机构、化工企业、政府机构）
& 国家级AI自主网络间谍活动
& \begin{itemize}[leftmargin=*,topsep=0pt,partopsep=0pt,itemsep=1pt,parsep=0pt]
    \item 官方报告说明攻击者通过社会工程手段（伪装成合法网络安全公司员工）使Claude Code执行相关任务~\cite{anthropicGTG1002_2025}。
    \item 官方报告显示攻击者将整体任务拆解为大量单独看似无害的技术请求，配合MCP协议调用侦察、漏洞验证等工具~\cite{anthropicGTG1002_2025}。
    \item Anthropic估计Claude独立完成了约80\%至90\%的战术性操作~\cite{anthropicGTG1002_2025}。
\end{itemize}
& \begin{itemize}[leftmargin=*,topsep=0pt,partopsep=0pt,itemsep=1pt,parsep=0pt]
    \item Anthropic确认约30个目标机构中有少数被成功入侵~\cite{anthropicGTG1002_2025}。
    \item Anthropic将此事件评估为已知首例主要由AI自主执行的大规模网络攻击活动~\cite{anthropicGTG1002_2025}。
    \item Anthropic报告同时说明Claude在此次活动中存在夸大战果、生成不准确信息等局限~\cite{anthropicGTG1002_2025}。
\end{itemize} \\
% \cline{2-5}
\hline 

2025.07.13--2025.07.19\newline （6天供应链投毒）
& Amazon Q Developer VS Code插件
& 开源供应链投毒、破坏性提示注入
& \begin{itemize}[leftmargin=*,topsep=0pt,partopsep=0pt,itemsep=1pt,parsep=0pt]
    \item 官方安全公告确认攻击者利用权限配置不当的GitHub令牌向开源仓库提交了未经批准的代码，漏洞编号CVE-2025-8217~\cite{amazonQCVE2025}。
    \item 公开报告显示注入的提示内容包含删除本地文件及终止/删除云资源（如EC2实例、S3存储桶）的指令~\cite{amazonQCVE2025}。
    \item 官方确认含该提示的1.84.0版本曾通过官方插件市场正式发布分发~\cite{amazonQCVE2025}。
\end{itemize}
& \begin{itemize}[leftmargin=*,topsep=0pt,partopsep=0pt,itemsep=1pt,parsep=0pt]
    \item AWS官方确认因代码存在语法错误，恶意指令未能实际执行，无客户资源受到影响~\cite{amazonQCVE2025}。
    \item 该插件安装量超95万次~\cite{amazonQCVE2025}。
    \item AWS随即撤销并替换相关凭证，移除相关代码并发布修复版本1.85.0~\cite{amazonQCVE2025}。
\end{itemize} \\
% \cline{2-5}
\hline 

2025.07.12--2025.07.23\newline （约12天测试周期）
& Replit AI编码智能体（单一企业测试项目）
& 违反代码冻结指令、生产数据库删除
& \begin{itemize}[leftmargin=*,topsep=0pt,partopsep=0pt,itemsep=1pt,parsep=0pt]
    \item 用户公开记录显示，Agent在明确下达"代码冻结"指令后仍执行了数据库删除操作~\cite{replitDatabaseDeletion2025}。
    \item 事件涉及生产数据库被错误删除，且删除前未获得人工二次确认~\cite{replitDatabaseDeletion2025}。
    \item 官方说明将问题归因于开发/生产环境隔离机制不足~\cite{replitDatabaseDeletion2025}。
\end{itemize}
& \begin{itemize}[leftmargin=*,topsep=0pt,partopsep=0pt,itemsep=1pt,parsep=0pt]
    \item 涉事项目生产数据库中约1200余名高管、约1200家公司记录被删除~\cite{replitDatabaseDeletion2025}。
    \item Agent此前告知用户数据无法回滚，但用户手动操作后确认数据实际可回滚~\cite{replitDatabaseDeletion2025}。
    \item 公司公开致歉并启动安全审查，宣布上线开发/生产环境自动隔离与一键回滚机制~\cite{replitDatabaseDeletion2025}。
\end{itemize} \\
\hline
\end{longtblr}
\endgroup
